% Appendix A: the formal names (loop0009 item 01). Every Lean-name footnote of
% the draft moved here; the text cites an entry as (Appendix A, A.n) with
% the leanref macro. check_lean.py and check_axioms.py read the leannames
% cells. Axiom line printed by paper2/axiom_check.lean (`make axioms`).
\section{Formal names}
\label{app:lean}

Every theorem, lemma and counterexample of this paper is checked in the Lean~4
proof assistant with the Mathlib library, unless the text says otherwise. The
reader need not read Lean to follow the paper. This appendix is for a reader
who wants to find the formal statement of a result: Table~\ref{tab:lean} maps
each result, in the order in which it appears, to the names of the Lean
declarations that state it. Where a row cites a definition rather than a
theorem, it is the definition the text uses.

All the theorems named depend only on the three standard axioms of Lean's
classical logic, \texttt{propext}, \texttt{Classical.choice} and
\texttt{Quot.sound}; an axiom check in the repository prints this for every
theorem named in the paper. The development has one \texttt{sorry}, Lean's marker for a statement
left unproved: a conjecture about lower bounds that no result of this paper uses. The part of
the development behind Section~\ref{sec:search} has no \texttt{sorry} at all.

Two Lean proofs take a route that the text does not describe. The proof of
Theorem~\ref{thm:decides} goes through narrowness \citep{LY2002ref11}, which
counts the closing order forwards rather than reversing it. The proof of
Theorem~\ref{thm:matching} reduces the deficiency form of Hall's theorem to
Mathlib's Hall theorem by adding $\delta$ dummy stacks. One Lean proof takes a
different route from the published one: the proof
that node search is monotone (Theorem~\ref{thm:core}(5)) follows the crusade
argument of \citet{BienstockSeymour1991}, not the reduction to edge search of
\citet{LaPaugh1993} on which \citet{LY2002ref10} rely. Three kinds of fact
are computed rather than proved. The values of $K_2$, $K_{1,3}$ and $K_{3,3}$
after Theorem~\ref{thm:bands}, which show that every end of the two bands
occurs, are computed by brute force from each quantity's own definition. The
search for smaller counterexamples (Appendix~\ref{app:minimality}) and the
results of Section~\ref{sec:results} are computations.

{\small
\begin{longtable}{@{}>{\raggedright}p{1.0cm}>{\raggedright}p{6.6cm}>{\raggedright\arraybackslash}p{8.0cm}@{}}
\caption{Every result of the paper has a formal statement in Lean, listed
here in the order in which the results appear.}\label{tab:lean}\\
\toprule
Entry & Result in the paper & Lean names \\
\midrule
\endfirsthead
\toprule
Entry & Result in the paper & Lean names \\
\midrule
\endhead
\bottomrule
\endfoot
\leangroup{Section~\ref{sec:equiv}, the equivalences}
\leanentry{clique}{Lemma~\ref{lem:clique}, cliques sit in a bag} \leannames{PathDecomposition.exists_bag_of_isClique} \\
\leanentry{vs}{Theorem~\ref{thm:vs}, vertex separation is pathwidth} \leannames{vertexSeparation_eq_pathwidth} \\
\leanentry{mosp}{Theorem~\ref{thm:mosp}, MOSP is pathwidth plus one} \leannames{MOSPInstance.mospValue_eq_pathwidth_add_one} \\
\leanentry{star}{The pattern graph gives no bound: the star $K_{1,4}$ (after Theorem~\ref{thm:mosp})} \leannames{star_mospValue, star_pathwidth_agreementGraph, star_refutes_pattern_graph_bound} \\
\leanentry{core-gml}{Theorem~\ref{thm:core}(1), gate matrix layout} \leannames{NetGateMatrix.tracks_eq_mospValue, NetGateMatrix.tracks_eq_pathwidth_add_one} \\
\leanentry{core-it}{Theorem~\ref{thm:core}(2), interval thickness} \leannames{intervalThickness_eq_pathwidth_add_one} \\
\leanentry{core-logic}{Theorem~\ref{thm:core}(3), one-dimensional logic} \leannames{LogicArray.tracks_eq_intervalThickness, LogicArray.tracks_eq_pathwidth_add_one} \\
\leanentry{core-narrow}{Theorem~\ref{thm:core}(4), narrowness} \leannames{narrowness_eq_pathwidth_add_one} \\
\leanentry{core-ns}{Theorem~\ref{thm:core}(5), node search} \leannames{nodeSearchMonotonicity, nodeSearch_eq_vertexSeparation_add_one, nodeSearch_chain} \\
\leanentry{core-fold}{Theorem~\ref{thm:core}(6), multiple folding} \leannames{NetGateMatrix.foldTracks_eq_pathwidth_add_one} \\
\leanentry{core-vsg}{Theorem~\ref{thm:core}(7), Lengauer's vertex separator game} \leannames{vsg_eq_max} \\
\leanentry{bands}{Theorem~\ref{thm:bands}, the bands} \leannames{pathwidth_le_splitBandwidth_le_pathwidth_add_one, vertexSeparation_le_edgeSearch_le_add_two, vertexSeparation_le_progressiveEdgeSearch_le_add_two} \\
\leanentry{false-pla}{Theorem~\ref{thm:false}(1), PLA folding} \leannames{plaTracks_idMatrix_five, plaTracks_idMatrix_unbounded, plaTracks_pathMatrix_unbounded} \\
\leanentry{false-cw}{Theorem~\ref{thm:false}(2), edge separation read as cutwidth} \leannames{cutwidth_unbounded, modCutwidth_unbounded} \\
\leangroup{Section~\ref{sec:search}, the search}
\leanentry{defs}{Step cost, $\Ext_k$ and $\Sol_k$ (Section~\ref{sec:search-def})} \leannames{stepCost, Solvable, SearchSol} \\
\leanentry{F}{Lemma~\ref{lem:F}, free moves never hurt} \leannames{solvable_iff_searchSol_cl} \\
\leanentry{decides}{Theorem~\ref{thm:decides}, the search decides pathwidth} \leannames{searchSol_empty_iff_pathwidth_add_one_le, searchSol_mospGraph_iff_mospValue_le} \\
\leanentry{dict1}{Table~\ref{tab:dictionary}: closing order from $\emptyset$} \leannames{orderOfLayout} \\
\leanentry{dict2}{Table~\ref{tab:dictionary}: open stacks $b(T)$} \leannames{openStacks_eq_card_boundary} \\
\leanentry{dict3}{Table~\ref{tab:dictionary}: $\cost(T,c)$} \leannames{stepCost_eq_card_boundary_add_one} \\
\leanentry{dict4}{Table~\ref{tab:dictionary}: $\Sol_k(\emptyset)$} \leannames{searchSol_empty_iff_pathwidth_add_one_le} \\
\leanentry{dict5}{Table~\ref{tab:dictionary}: free closure $\cl(T)$} \leannames{mem_cl_iff} \\
\leanentry{dict6}{Table~\ref{tab:dictionary}: definite move, repaired} \leannames{isHereditarilyDefinite_iff_isCommittable} \\
\leanentry{premise}{Lemma~\ref{lem:premise}, what the premise says} \leannames{closeCount_eq, isDefinite_iff} \\
\leanentry{cex}{Counterexample~\ref{cex:definite}} \leannames{definiteMove_counterexample, cexFamily_closed} \\
\leanentry{thm1false}{Chu and Stuckey's Theorem~1 (the definite move), under both readings, and Chu's Theorem~6.3.6 are false} \leannames{chuStuckey_theorem1_false, chuStuckey_theorem1_false_literal, chuThesis_theorem636_false} \\
\leanentry{fink}{Fink's restatement of the definite move is false} \leannames{fink_theorem1_false} \\
\leanentry{hered}{The hereditarily definite premise, and its failure in Counterexample~\ref{cex:definite}} \leannames{IsHereditarilyDefinite, not_isHereditarilyDefinite_cex} \\
\leanentry{submod}{Lemma~\ref{lem:submod}, open stacks are submodular} \leannames{openStacks_submodular} \\
\leanentry{repaired}{Theorem~\ref{thm:repaired}, the repaired definite move is sound} \leannames{solvable_cl_insert_of_hereditarilyDefinite, searchSol_cl_insert_of_hereditarilyDefinite} \\
\leanentry{matching}{Theorem~\ref{thm:matching}, the repair is a matching condition} \leannames{isHereditarilyDefinite_iff_hasDefiniteMatching} \\
\leanentry{open1}{The published definite move is right when $\open(q,S)\le1$} \leannames{isHereditarilyDefinite_of_openCount_le_one} \\
\leanentry{subset}{The subset rule is sound} \leannames{searchSol_cl_insert_of_newlyOpened_subset, subsetFilter_sound} \\
\leanentry{tiebreak}{The subset rule needs its tie-break (7-customer graph)} \leannames{noTieBreak_counterexample} \\
\leanentry{isbetter}{The better move's premise, restated} \leannames{isBetter_iff} \\
\leanentry{bettercex}{Chu and Stuckey's Theorem~2 (the better move) is false} \leannames{betterMove_counterexample, chuStuckey_theorem2_false} \\
\leanentry{thm638}{Chu's Theorem~6.3.8 is false} \leannames{chuThesis_theorem638_false, chuThesis_theorem638_false_literal} \\
\leanentry{repbetter}{The repaired better move is sound} \leannames{IsRepairedBetter, searchSol_cl_insert_of_repairedBetter} \\
\leanentry{bugs}{The better move's two requirements: the close count and the rule order} \leannames{bugA_counterexample, bugB_counterexample} \\
\leanentry{oldmove}{The old move is sound} \leannames{searchSol_reinsert, not_searchSol_of_oldMove} \\
\leanentry{reinsert}{The old move needs its reinsertion test (path on five vertices)} \leannames{reinsert_needs_test} \\
\leanentry{memo}{Solvability of a state does not depend on the path to it (the memo)} \leannames{solvable_iff_of_cl_eq} \\
\leanentry{fakeold}{A node with an unrefuted old move can record a solvable state} \leannames{exec_fake_oldMove} \\
\leanentry{nodesound}{Node-soundness of a filter (Section~\ref{sec:theorem})} \leannames{NodeSoundAt} \\
\leanentry{main}{Theorem~\ref{thm:main}, the repaired search is sound} \leannames{exec_repairedFullFilter_sound, exec_repairedFullFilter_pathwidth, exec_repairedFullFilter_mospValue} \\
\leanentry{notsound}{Under the published rules the filter is not node-sound} \leannames{not_codeFilterSound_cexGraph} \\
\leanentry{endpoint}{Lemma~\ref{lem:endpoint}, the premises as commitment conditions} \leannames{isDefinite_iff_endpoint, isHereditarilyDefinite_iff_isCommittable} \\
\leanentry{cexnotcommit}{In Counterexample~\ref{cex:definite} the endpoint test passes and the commitment fails} \leannames{cex_isDefinite_not_isCommittable} \\
\leanentry{existscommit}{The published premise gives a commitment to the least-border set} \leannames{exists_isCommittable_of_isDefinite} \\
\leanentry{commit234}{In Counterexample~\ref{cex:definite} that set is $\{2,3,4\}$} \leannames{cex_isCommittable_234} \\
\leanentry{depth1}{A commitment of depth one holds iff $\open(q,S)\le1$} \leannames{isCommittable_insert_iff} \\
\leanentry{greedy}{The greedy step of Coudert, Mazauric and Nisse is that condition} \leannames{isGreedyStep_iff} \\
\leangroup{Section~\ref{sec:results}, computational results}
\leanentry{split}{The split of the root into independent tasks is sound} \leannames{execSplit_repairedFullFilter_sound, execSplit_repairedFullFilter_mospValue} \\
\leangroup{Appendix~\ref{app:machinery-all}, supporting material}
\leanentry{edgeless-ns}{Node search equalities need an edge (Appendix~\ref{app:edge})} \leannames{nodeSearch_ne_intervalThickness_of_edgeless, monotoneNodeSearch_ne_vertexSeparation_add_one_of_edgeless} \\
\leanentry{vsg0}{Lengauer's Theorem~4 fails at $k=0$ (Appendix~\ref{app:edge})} \leannames{isPositiveVSG_triangleGraph_counterexample} \\
\leanentry{tns}{$t=\ns$ fails on the matrix $[1]$ (Appendix~\ref{app:edge})} \leannames{tracks_ne_monotoneNodeSearch_one} \\
\leanentry{kp}{The Kirousis--Papadimitriou claim fails on $K_{1,3}$ (Appendix~\ref{app:edge})} \leannames{kirousisPapadimitriou_claim2_false} \\
\leanentry{pebbling}{Progressive pebbling is in the exact core (Appendix~\ref{app:thirteenth})} \leannames{pbw_lengauerD_eq_pathwidth, pbw_eq_mospValue_pebbleMatrix, mpb_eq_pathwidth_add_one} \\
\leanentry{wavefront}{The maximum wavefront is in the exact core (Appendix~\ref{app:thirteenth})} \leannames{card_wavefront, maxWavefront_eq_vertexSepOfLayout_add_one, minMaxWavefront_eq_pathwidth_add_one, minMaxWavefront_eq_narrowness, minMaxWavefront_of_isEmpty} \\
\leanentry{filter}{Proposition~\ref{prop:filter}, the repaired filter is node-sound} \leannames{repairedFullFilter_sound} \\
\leanentry{runs}{Theorem~\ref{thm:runs}, runs are sound} \leannames{Exec.sound} \\
\leanentry{encoding}{Appendix~\ref{app:sat}, the SAT formulation is faithful (clauses 1--4)} \leannames{encodes_iff_mospValue_le} \\
\end{longtable}
}
